% Cumulative TRIT extension; inserted before the regime figure.
% It neither replaces nor removes the preceding development.
\subsubsection{Balanced division and composition of lifts}
\label{trit-ampl-div-balanceada}

Ternary reduction expresses an integer evaluation of APP through a visible
coordinate and a prolongation coordinate. Its function is not merely to
abbreviate the notation for an integer: it provides a composition rule
determining how the quotient changes when evaluations are added or multiplied.
The balanced choice also makes the compatibility of this rule with sign
reversal explicit.

\begin{definition}[Balanced ternary coordinates]
\label{trit-ampl-def-coordenadas}
For \(n\in\mathbb Z\), let
\[
 q_3(n)=\left\lfloor\frac{n+1}{3}\right\rfloor,\qquad
 r_3(n)=n-3q_3(n).
\]
The map
\[
 \mathcal B:\mathbb Z\longrightarrow\TRIT\times\mathbb Z,\qquad
 n\longmapsto(r_3(n),q_3(n))
\]
is called the balanced ternary decomposition. Its inverse is
\(\mathcal L(r,c)=r+3c\).
\end{definition}

Indeed, \(r_3(n)\in\{-1,0,+1\}\). If two pairs represent the same integer,
\(r+3c=s+3d\), then \(r-s\) is a multiple of \(3\) and belongs to the interval
\([-2,2]\); necessarily \(r=s\) and \(c=d\). The representation is therefore
unique. Division can be repeated on \(q_3(n)\), so that each level preserves
the data needed to reconstruct the preceding one. The quotient is not a
perturbation of the residue: it is its complementary coordinate.

The signature obtained from an evaluation \(F\), with \(F=S\) or \(F=P\) in APP,
is \(r_3\circ F\). When reduction acts on the phase cursor, the same balanced
section produces the classes \(\{1,4,7\}\), \(\{2,5,8\}\), and
\(\{3,6,9\}\), with values \(+1,-1,0\). The TPK selector defined below
uses these classes to determine the operating mode.
Reading an evaluation and selecting a mode thus share the same ternary
reduction, but have different domains and functions.

\Needspace{16\baselineskip}
\begin{proposition}[Arithmetic of balanced pairs]
\label{trit-ampl-prop-aritmetica}
Let \(n=r+3c\) and \(m=s+3d\), with \(r,s\in\TRIT\). Define
\[
 \chi(r,s)=\frac{r+s-r_3(r+s)}3.
\]
Integer operations are transported to the pairs by
\begin{align}
 (r,c)\boxplus(s,d)
 &=\bigl(r_3(r+s),c+d+\chi(r,s)\bigr),
 \label{trit-ampl-eq-suma}\\
 (r,c)\boxtimes(s,d)
 &=\bigl(rs,rd+sc+3cd\bigr).
 \label{trit-ampl-eq-producto}
\end{align}
The map \(\mathcal L\) preserves both operations.
\end{proposition}
\begin{proof}
The equality \(r+s=r_3(r+s)+3\chi(r,s)\) gives
\[
 n+m=r_3(r+s)+3\bigl(c+d+\chi(r,s)\bigr).
\]
On the other hand,
\[
 nm=rs+3(rd+sc+3cd).
\]
Since the product of two balanced representatives again belongs to
\(\{-1,0,+1\}\), it requires no additional residue reduction. Uniqueness
of the decomposition proves both formulas.
\end{proof}

For example, \(2=(-1)+3\) and \(4=1+3\). Their sum and product are reconstructed
as
\[
 (-1,1)\boxplus(1,1)=(0,2),\qquad
 (-1,1)\boxtimes(1,1)=(-1,3).
\]
The resulting evaluations are \(6\) and \(8\). Adding the quotients suffices
in this additive example; the product involves the cross terms
\(rd\), \(sc\), and \(3cd\). The discrepancy between the two operations
also resides in the prolongation coordinate, not only in the residue.

The term \(\chi\) is a cocycle of the balanced section. If
\(r\oplus s=r_3(r+s)\), it satisfies
\[
 \chi(r,s)+\chi(r\oplus s,t)
 =\chi(s,t)+\chi(r,s\oplus t).
\]
Both sides represent the quotient of the same sum of three integers.
This identity ensures that regrouping the operations does not change the
reconstructed evaluation. At the same time, projection onto the first
component suppresses this accounting: \(1+1=-1+3\), whereas the residue
reading retains only \(-1\). The information removed is determined,
in this case, by \(\chi(1,1)=1\).

The relation with modulus \(9\) preserves a second discrete scale.
The map
\[
 \beta:\TRIT\times\mathbb F_3\longrightarrow\mathbb Z/9\mathbb Z,
 \qquad (r,\bar c)\longmapsto r+3c\pmod9
\]
is well defined and bijective. Changing the representative of \(\bar c\)
adds a multiple of \(9\); conversely, equality of images fixes
\(r\) first, by reduction modulo \(3\), and then \(\bar c\). Addition
transported by this bijection contains the term
\(\chi(r,s)\) from \eqref{trit-ampl-eq-suma}, now reduced modulo \(3\).
Thus the two ternary levels do not compose by independent coordinatewise
addition: normalization of the first level modifies the second.

In particular, the classes \(3\), \(6\), and \(0\) of
\(\mathbb Z/9\mathbb Z\) have first component \(0\), but second components
\(1\), \(2\), and \(0\), respectively. The distinction already observed
between direct reduction and extraction of the ideal coordinate appears
as a property of \(\beta\). Preserving this coordinate prevents the zero
residue sector from being artificially collapsed into a single position.
Modulus \(1000\), used later to organize decimal blocks,
serves a different positional-representation function; its quotient is
preserved through the corresponding division. A change of base preserves
the integer when it retains residue and quotient, whereas coincidence of
residues in different bases does not by itself identify their transport states.

\Needspace{11\baselineskip}
\subsubsection{Inversion, residue class, and non-visible coordinates}
\label{trit-ampl-inversion}

Additive inversion lifts unambiguously:
\[
 \mathcal B(-n)=(-r_3(n),-q_3(n)).
\]
In particular, \(3\), \(0\), and \(-3\) have coordinates
\((0,1)\), \((0,0)\), and \((0,-1)\). They share a zero signature but
represent different evaluations. In a transition restricted to a quotient
of maximum amplitude \(1\), these are the three possibilities of the
visible fiber over zero; in a general integer lift, that fiber contains
all pairs \((0,c)\), with \(c\in\mathbb Z\).

The quotient coordinate recovers the evaluated integer, but does not by
itself recover the trajectory that produced it. Two APP paths may share
evaluation, residue, and quotient while differing in orientation, sequence
of modes, or winding. Thus decomposition \(\mathcal B\) is integrated into
the transport state rather than replacing it. The arithmetic memory of a
division and the memory of a trajectory are related coordinates, not synonyms.

For a signature sequence of length \(n\), oriented reversal is
\[
 C_{\rm rev}(\tau_1,\ldots,\tau_n)
   =(-\tau_n,\ldots,-\tau_1).
\]
Applying it twice restores the sequence. The operation reverses the order
of positions and the sign of each signature; it preserves the length and
the zero positions in reversed order. In the adopted parametrization,
it interchanges return with memory, associated with \(+1\), and bidirectional
opening, associated with \(-1\), while retaining the threshold \(0\).
Transport of the remaining trajectory data is specified by the corresponding
TPK involution.

This signature reversal must be distinguished from conjugation within a
quadratic algebra. The former may change the type \(\tau\); the latter,
constructed below, acts within a type already fixed. This distinction
prevents attributing an algebraic equivalence between different regimes
to a sign reversal.

\subsubsection{Quadratic product, conjugation, and algebraic norm}
\label{trit-ampl-norma}

An element of \(\mathbb A_\tau\) has a unique expression \(aI+bJ_\tau\).
The relation \(J_\tau^2=-\tau I\) determines its product:
\begin{equation}
 (aI+bJ_\tau)(cI+dJ_\tau)
   =(ac-\tau bd)I+(ad+bc)J_\tau.
 \label{trit-ampl-producto-cuadratico}
\end{equation}
This is a real realization of the type already selected. The coordinates
\(a,b\) are not APP residues, and the product
\eqref{trit-ampl-producto-cuadratico} does not replace
\eqref{trit-ampl-eq-producto}; it expresses another operation in the
declared quadratic codomain.

\begin{definition}[Conjugation and algebraic norm]
\label{trit-ampl-def-norma}
Conjugation of type \(\tau\) and its algebraic norm are
\[
 \overline{aI+bJ_\tau}=aI-bJ_\tau,\qquad
 N_\tau(aI+bJ_\tau)=a^2+\tau b^2.
\]
\end{definition}

\begin{proposition}[Multiplicativity and invertibility]
\label{trit-ampl-prop-norma}
For \(z,w\in\mathbb A_\tau\),
\[
 z\bar z=N_\tau(z)I,\qquad
 N_\tau(zw)=N_\tau(z)N_\tau(w).
\]
Moreover, \(z\) is invertible if and only if \(N_\tau(z)\ne0\), in which case
\[
 z^{-1}=\frac{\bar z}{N_\tau(z)}.
\]
\end{proposition}
\begin{proof}
The product of \(aI+bJ_\tau\) and \(aI-bJ_\tau\) is
\((a^2+\tau b^2)I\). Conjugation is multiplicative and the algebra is
commutative, so that
\((zw)\overline{zw}=(z\bar z)(w\bar w)\). If the norm does not vanish,
the stated formula supplies the inverse. Conversely, if \(zw=I\),
multiplicativity implies \(N_\tau(z)N_\tau(w)=1\).
\end{proof}

The quadratic form \(N_\tau\) is positive definite only in the elliptic type.
For \(\tau=0\), it equals \(a^2\), and the element \(J_0\) is nonzero although
its square and norm are zero. For \(\tau=-1\), it equals \(a^2-b^2\);
the nonzero elements \(I+J_{-1}\) and \(I-J_{-1}\) have zero norm and
their product vanishes. Here the algebraic norm is a multiplicative
quadratic form: it does not presuppose a positive vector-space norm.

The three types admit faithful realizations by
\[
 aI+bJ_\tau\longmapsto
 \begin{pmatrix}a&-\tau b\\b&a\end{pmatrix}.
\]
The determinant of this matrix is \(N_\tau(aI+bJ_\tau)\).
The complex structure, the dual-number structure, and the split real
structure are thus represented by a single matrix expression.
Each preserves its own law and zero divisors. The form of signature
\((1,1)\) in the last case admits a Lorentzian interpretation in dimension
\(1+1\); identification of a physical metric additionally requires a
specific space, field, and transport rule.

\subsubsection{Multiplicative conservation and composition of evolutions}
\label{trit-ampl-evoluciones}

The exponential already obtained belongs to the set of norm-one elements:
\[
 N_\tau\bigl(\exp(tJ_\tau)\bigr)=1.
\]
This conservation persists under composition of parameters. Indeed, all
factors are functions of the same generator, and
\[
 \exp(tJ_\tau)\exp(sJ_\tau)=\exp((t+s)J_\tau).
\]
Comparing the two components gives the addition laws
\begin{align}
 C_\tau(t+s)&=C_\tau(t)C_\tau(s)-\tau S_\tau(t)S_\tau(s),\\
 S_\tau(t+s)&=S_\tau(t)C_\tau(s)+C_\tau(t)S_\tau(s).
\end{align}
Conjugation changes \(t\) to \(-t\) and coincides with inversion within
this subgroup. It does not change \(\tau\): it reverses an evolution of
the chosen regime.

In the elliptic type, the preserved form is \(a^2+b^2\), and the exponential
trajectory traverses the unit-norm circle. In the hyperbolic type, the
eigencoordinates are \(u=a+b\), \(v=a-b\), with \(uv=N_{-1}(z)\).
The exponential gives
\[
 (u,v)=(e^t,e^{-t}),\qquad uv=1.
\]
Growth of one coordinate is accompanied by contraction of the other.
Conservation expresses a relation between the two coordinates; it does
not assert that each remains constant.

In the parabolic type,
\[
 (I+tJ_0)(I+sJ_0)=I+(t+s)J_0,\qquad
 (I+tJ_0)^{-1}=I-tJ_0.
\]
The evolution does not stop because \(J_0^2=0\). Nilpotency removes terms
of order at least two, but preserves the linear term that transports
the parameter. The visible symbol \(0\), the zero element of the algebra,
and the nonzero generator \(J_0\) are distinct. This distinction is
necessary to describe the threshold without identifying it with an absence
of state or dynamics.

\begin{proposition}[Composition of affine coordinates]
\label{trit-ampl-prop-pendientes}
In the chart where the scalar component does not vanish, represent a
direction by \(I+uJ_\tau\). If \(1-\tau uv\ne0\), the direction of its
product with \(I+vJ_\tau\) has coordinate
\[
 u\star_\tau v=\frac{u+v}{1-\tau uv}.
\]
\end{proposition}
\begin{proof}
We have
\[
 (I+uJ_\tau)(I+vJ_\tau)
   =(1-\tau uv)I+(u+v)J_\tau.
\]
Dividing by the scalar component gives the stated expression.
\end{proof}

The affine chart preserves the ratio of components, not the scalar factor
divided out. When \(1-\tau uv=0\), the product must be examined before
changing charts: it may have a well-defined direction with zero scalar
component, or be the zero element in the presence of zero divisors.
The composition law still belongs to the full algebra; a coordinate
singularity does not automatically imply a singularity of the object.
This formula will reappear in the elliptic regional reading. Its use
requires retaining the denominator and, where applicable, the normalization factor.

\subsubsection{Transport of type and subsequent realizations}
\label{trit-ampl-puentes}

The algebras \(\mathbb A_{+1}\), \(\mathbb A_0\), and
\(\mathbb A_{-1}\) are not isomorphic as unital real algebras.
The first is a field; the second contains a nonzero nilpotent; the
third has nontrivial idempotents and is isomorphic to
\(\mathbb R\times\mathbb R\). Thus a change of signature organizing a
route cannot be interpreted as an undifferentiated isomorphism between
the three realizations. The TPK preserves the regime index, the domains,
and the effective operator acting at each transition.

In the electronic development of Article I of this series, the pair of projections of the APP sheets
produces a normalized commutator whose square is \(-I\). The same block
contains an involution of square \(I\) and two nilpotent directions.
The joint realization is proved there for that concrete pair of projections;
the preceding classification explains why the three quadratic relations
can appear within one matrix algebra without being identified with one another.

In the exceptional prolongation, the Paley matrix reduced modulo \(3\)
satisfies \(A_W^2=-I_6\). This relation endows the ternary space with
a structure of dimension \(3\) over \(\mathbb F_9\). Its connection to
a real operator root is expressed through the integral presentation
\[
 \mathbb Z[u]/(u^2+1).
\]
Base changes to \(\mathbb F_3\) and \(\mathbb R\) produce representations
of the same relation. Each retains its characteristic, dimension, and
incidence information. Developing the respective operators establishes
the connection, not equality of their cardinalities.

Norm-one conservation consequently describes the quadratic realizations
of evolutions. Conservation of quotients, orientation, and trajectory
belongs to the transported arithmetic state. Both properties participate
in the construction, with different functions: one determines the identities
of the representations; the other makes it possible to compose, reconstruct,
and prolong the operations that give rise to them.
